\section{路径 C：消融矩阵与统一实验矩阵（P1，纯本地计算）}

\subsection{C1：组件消融跨数据集（改动量最小、收益直接）}

现状：M3（\texttt{run\_journal\_analysis.py}）把组件消融硬编码在 QASPER 池。改动：把消融配置从``QASPER only''参数化为三池循环（配置清单不变：全配置、$\alpha{=}0$ 去问题条件化、floor=1 平位置先验、$\beta{=}0$ 去冗余、纯相关性 $\alpha{=}1$+平先验+$\beta{=}0$）。改动量约为该脚本内一个循环边界与结果分组的十几行；池文件、选择器实现、指标计算全部复用。产出：三数据集 $\times$ 五配置的消融表（每格证据召回加 bootstrap 区间）。预期与写作要点：HotpotQA 与 2Wiki 的问题实体命名更强，问题条件化的消融差应更大；位置先验在两 wiki 池上的``召回代价''可能不同于 QASPER（证据分布位置不同，M6 的池统计会给出预期方向）。无论方向如何，跨数据集一致的消融结构直接支撑``三目标是普遍机制而非单数据集偶然''的声明。

\subsection{C2：配置 \texorpdfstring{$\times$}{x} 预算交叉矩阵}

现状：M2 预算扫描只覆盖五种选择器（sieve/head/textrank/random/bm25）在 $r\in\{2,4,8,16\}$；M3 消融只在 $r{=}4$。改动：把消融五配置加入预算扫描的循环体，形成``配置（5 消融 + 5 基线 + GPT-2 门）$\times$ 预算（4 档）$\times$ 数据集（3 池）''的完整召回矩阵；GPT-2 门复用已落盘的 100 条测量（$r{=}4$ 单点），不重跑。纯本地计算的总代价估算：每格一次选择器运行（毫秒级）加 bootstrap（1{,}000 重采样，秒级），全矩阵约 200 个``配置 $\times$ 预算 $\times$ 数据集''格，单容器数小时内完成。产出论文新表（或扩表现表）：每行一个配置，每列一个预算，报告召回与配对区间；配套把时序测量（\texttt{run\_timing\_journal.py} 的既有数字）并入行属性，使每行天然携带``召回--压缩--成本''三元组——这正是三轴 frontier 图的数据源，C2 因此同时是效率节的补强。

\subsection{C3：frontier 图升级}

论文现有的 frontier 图（fig4）只画了少数点。C2 矩阵落盘后重绘：横轴选择成本（对数刻度，CPU ms 到 GPT-2 秒），纵轴证据召回（或端到端 F1 保留率，若 A 路径完成），同一预算的配置连成等预算线，三个数据集三张子图。预期图景：相关性族（bm25、纯相关性消融、sieve 全配置的前半段）聚在高召回带，全配置在等预算线上向``放置与多样性''方向位移，神经门在成本轴右端——三根支柱（formulation、frontier、结构性权衡）在同一张图上各自获得视觉证据。图重绘用 \texttt{fig\_journal.py} 既有管线，改数据源即可。

\subsection{C4：可选——预算自适应实验（formulation 的自然延伸）}

第一阶段审计确认``自适应预算''已被 AdaComp 占位（文档级、检索质量驱动）。若时间允许，可在 \sieve{} 框架内做词预算随问题复杂度自适应的最小实验（例如按问题 token 数或查询词 IDF 和分桶调 $B$），作为 formulation 的延伸而非新颖性主张——它证明三目标框架可以吸收预算维度，而与 AdaComp 的差异（句子级三目标 vs 文档级自适应）在 related work 已划定。此项不影响 DoD，仅在有富余算力时执行。
